% META: {"id": "1NSI-TC-EX-055", "chapitre": "1NSI-TYPES-CONSTRUITS", "type_objet": "exercice", "capacites": ["1NSI-TYPES-CONSTRUITS-C5"], "parcours": 3, "competences": ["programmer", "modeliser"], "duree_min": 20, "mode_creation": "ex_nihilo", "sources_inspiration": [], "status": "needs_review"}
% BEGIN-VERIFY
% def historique_scores(equipes, resultats):
%     # equipes est une liste de dicos avec cle "nom" et "scores" (liste)
%     # resultats est une liste de tuples (nom, score)
%     # on ajoute chaque score a la bonne equipe, sans effet de bord
%     copie = []
%     for e in equipes:
%         copie.append({"nom": e["nom"], "scores": list(e["scores"])})
%     for nom, score in resultats:
%         for e in copie:
%             if e["nom"] == nom:
%                 e["scores"].append(score)
%     return copie
% equipes = [{"nom": "A", "scores": [10, 20]},
%             {"nom": "B", "scores": [15]}]
% nouvelles = historique_scores(equipes, [("A", 30), ("B", 25)])
% assert equipes[0]["scores"] == [10, 20]
% assert nouvelles[0]["scores"] == [10, 20, 30]
% assert equipes[1]["scores"] == [15]
% assert nouvelles[1]["scores"] == [15, 25]
% END-VERIFY
\begin{exercice}{1NSI-TC-EX-055}{3}{20}
On représente des équipes d'e-sport par une liste de dictionnaires ayant les clés
\lstinline{"nom"} et \lstinline{"scores"} (une liste d'entiers).

\begin{enumerate}
  \item Écrire une fonction \lstinline{historique_scores(equipes, resultats)} qui :
        \begin{itemize}
          \item prend la liste des équipes et une liste de tuples
                \lstinline{(nom, score)} ;
          \item renvoie une \textbf{nouvelle} liste d'équipes où les scores ont été
                ajoutés ;
          \item \textbf{ne modifie pas} la liste d'origine (ni les dictionnaires,
                ni les listes de scores qu'ils contiennent).
        \end{itemize}
  \item Tester avec :
\begin{python}
equipes = [{"nom": "A", "scores": [10, 20]},
            {"nom": "B", "scores": [15]}]
nouvelles = historique_scores(equipes, [("A", 30), ("B", 25)])
print(equipes[0]["scores"])   # [10, 20]
print(nouvelles[0]["scores"]) # [10, 20, 30]
\end{python}
  \item Expliquer pourquoi il faut copier à la fois les dictionnaires et les listes
        de scores.
\end{enumerate}
\end{exercice}
